% Options for packages loaded elsewhere
\PassOptionsToPackage{unicode}{hyperref}
\PassOptionsToPackage{hyphens}{url}
\PassOptionsToPackage{dvipsnames,svgnames,x11names}{xcolor}
\documentclass[
  10pt,
  fontset=fandol]{ctexart}
\usepackage{xcolor}
\usepackage[margin=16mm]{geometry}
\usepackage{amsmath,amssymb}
\setcounter{secnumdepth}{-\maxdimen} % remove section numbering
\usepackage{iftex}
\ifPDFTeX
  \usepackage[T1]{fontenc}
  \usepackage[utf8]{inputenc}
  \usepackage{textcomp} % provide euro and other symbols
\else % if luatex or xetex
  \usepackage{unicode-math} % this also loads fontspec
  \defaultfontfeatures{Scale=MatchLowercase}
  \defaultfontfeatures[\rmfamily]{Ligatures=TeX,Scale=1}
\fi
\usepackage{lmodern}
\ifPDFTeX\else
  % xetex/luatex font selection
\fi
% Use upquote if available, for straight quotes in verbatim environments
\IfFileExists{upquote.sty}{\usepackage{upquote}}{}
\IfFileExists{microtype.sty}{% use microtype if available
  \usepackage[]{microtype}
  \UseMicrotypeSet[protrusion]{basicmath} % disable protrusion for tt fonts
}{}
\makeatletter
\@ifundefined{KOMAClassName}{% if non-KOMA class
  \IfFileExists{parskip.sty}{%
    \usepackage{parskip}
  }{% else
    \setlength{\parindent}{0pt}
    \setlength{\parskip}{6pt plus 2pt minus 1pt}}
}{% if KOMA class
  \KOMAoptions{parskip=half}}
\makeatother
\setlength{\emergencystretch}{3em} % prevent overfull lines
\providecommand{\tightlist}{%
  \setlength{\itemsep}{0pt}\setlength{\parskip}{0pt}}
\usepackage{amsmath,amssymb,mathtools,needspace}
\xeCJKDeclareCharClass{CJK}{"2032,"0400->"04FF,"2160->"216B,"2460->"2473,"25A0,"25C7,"25CB,"25CF}
\IfFontExistsTF{Arial Unicode MS}{\xeCJKsetup{AutoFallBack=true}\setCJKfallbackfamilyfont{\CJKrmdefault}{Arial Unicode MS}}{}
\setcounter{secnumdepth}{0}
\setlength{\emergencystretch}{2em}
\allowdisplaybreaks
\usepackage{bookmark}
\IfFileExists{xurl.sty}{\usepackage{xurl}}{} % add URL line breaks if available
\urlstyle{same}
\hypersetup{
  pdftitle={消元手续},
  colorlinks=true,
  linkcolor={Maroon},
  filecolor={Maroon},
  citecolor={Blue},
  urlcolor={Blue},
  pdfcreator={LaTeX via pandoc}}

\title{消元手续}
\author{}
\date{}

\begin{document}
\maketitle

\subsubsection{7.2.1 消元手续}\label{ux6d88ux5143ux624bux7eed}

设有线性方程组

\[
\begin{cases}
a_{11}x_1+a_{12}x_2+\cdots+a_{1n}x_n=b_1,\\
a_{21}x_1+a_{22}x_2+\cdots+a_{2n}x_n=b_2,\\
\qquad\qquad\vdots\\
a_{n1}x_1+a_{n2}x_2+\cdots+a_{nn}x_n=b_n,
\end{cases}
\tag{7.2.1}
\]

或写成矩阵形式 \(Ax=b\)，其中

\[
A=\begin{pmatrix}
a_{11}&a_{12}&\cdots&a_{1n}\\
a_{21}&a_{22}&\cdots&a_{2n}\\
\vdots&\vdots&&\vdots\\
a_{n1}&a_{n2}&\cdots&a_{nn}
\end{pmatrix},\quad
x=\begin{pmatrix}x_1\\x_2\\\vdots\\x_n\end{pmatrix},\quad
b=\begin{pmatrix}b_1\\b_2\\\vdots\\b_n\end{pmatrix},
\]

\(A\) 为非奇异矩阵。下面举一个简单的例子来说明消去法的基本思想。

\textbf{例 7.1} 用消去法解方程组

\[x_1+x_2+x_3=6,\tag{7.2.2}\]

\[4x_2-x_3=5,\tag{7.2.3}\]

\[2x_1-2x_2+x_3=1.\tag{7.2.4}\]

\textbf{解} 第一步，将式（7.2.2）乘以 \(-2\) 加到式（7.2.4）上去，消去式（7.2.4）中的未知数 \(x_1\)，得到

\[-4x_2-x_3=-11.\tag{7.2.5}\]

第二步，将式（7.2.3）加到式（7.2.5）上，消去式（7.2.5）中的未知数 \(x_2\)，得到与原方程组等价的三角方程组

\[
\begin{cases}
x_1+x_2+x_3=6,\\
4x_2-x_3=5,\\
-2x_3=-6.
\end{cases}
\tag{7.2.6}
\]

显然方程组（7.2.6）是容易求解的，解为 \(x^*=(1,2,3)^{\mathrm T}\)。上述过程相当于

\[
(A\mathbin{\vdots}b)=
\left(\begin{array}{rrr|r}1&1&1&6\\0&4&-1&5\\2&-2&1&1\end{array}\right)
\longrightarrow
\left(\begin{array}{rrr|r}1&1&1&6\\0&4&-1&5\\0&-4&-1&-11\end{array}\right)
\longrightarrow
\left(\begin{array}{rrr|r}1&1&1&6\\0&4&-1&5\\0&0&-2&-6\end{array}\right).
\]

\[(-2)\times r_1{}^{\text{①}}+r_3\to r_3,\qquad r_2+r_3\to r_3.\]

由此看出，用消去法解方程组的基本思想是，用逐次消去未知数的方法把原来方程组 \(Ax=b\) 化为与其等价的三角方程组，而求解三角方程组就容易了。换句话说，上

\begin{quote}
① \(r_i\) 表示矩阵的第 \(i\) 行。
\end{quote}

\href{../../source-images/pdf-179.jpeg}{原页扫描}

述过程就是用行的初等变换将原方程组系数矩阵化为简单形式，从而将求解原方程组（7.2.1）的问题转化为求解简单方程组的问题。

下面来讨论一般的解 \(n\) 阶方程组的 Gauss 消去法。

将式（7.2.1）记作 \(A^{(1)}x=b^{(1)}\)，其中 \(A^{(1)}=(a_{ij}^{(1)})=(a_{ij})\)，\(b^{(1)}=b\)。

（1）第一次消元。设 \(a_{11}^{(1)}\ne0\)，首先对行计算乘数 \(m_{i1}=a_{i1}^{(1)}/a_{11}^{(1)}\)（\(i=2,3,\cdots,n\)），用 \(-m_{i1}\) 乘式（7.2.1）的第 1 个方程，加到第 \(i\)（\(i=2,3,\cdots,n\)）个方程上，消去式（7.2.1）的第 2 个方程直到第 \(n\) 个方程中的未知数 \(x_1\)，得与式（7.2.1）等价的方程组

\[
\begin{pmatrix}
a_{11}^{(1)}&a_{12}^{(1)}&\cdots&a_{1n}^{(1)}\\
0&a_{22}^{(2)}&\cdots&a_{2n}^{(2)}\\
\vdots&\vdots&&\vdots\\
0&a_{n2}^{(2)}&\cdots&a_{nn}^{(2)}
\end{pmatrix}
\begin{pmatrix}x_1\\x_2\\\vdots\\x_n\end{pmatrix}
=\begin{pmatrix}b_1^{(1)}\\b_2^{(2)}\\\vdots\\b_n^{(2)}\end{pmatrix},
\tag{7.2.7}
\]

简记作 \(A^{(2)}x=b^{(2)}\)，其中

\[
a_{ij}^{(2)}=a_{ij}^{(1)}-m_{i1}a_{1j}^{(1)},\quad
b_i^{(2)}=b_i^{(1)}-m_{i1}b_1^{(1)}\quad(i,j=2,3,\cdots,n).
\]

（2）一般第 \(k\)（\(1\le k\le n-1\)）次消元。设第 \(k-1\) 步计算已经完成，即已计算好与式（7.2.1）等价的方程组

\[A^{(k)}x=b^{(k)},\tag{7.2.8}\]

且已消去未知数 \(x_1,x_2,\cdots,x_{k-1}\)，其中 \(A^{(k)}\) 具有如下形式：

\[
A^{(k)}=\begin{pmatrix}
a_{11}^{(1)}&a_{12}^{(1)}&\cdots&\cdots&\cdots&a_{1n}^{(1)}\\
&a_{22}^{(2)}&\cdots&\cdots&\cdots&a_{2n}^{(2)}\\
&&\ddots&&&\vdots\\
&&&a_{kk}^{(k)}&\cdots&a_{kn}^{(k)}\\
&&&\vdots&&\vdots\\
&&&a_{nk}^{(k)}&\cdots&a_{nn}^{(k)}
\end{pmatrix}.
\]

设 \(a_{kk}^{(k)}\ne0\)，计算乘数 \(m_{ik}=a_{ik}^{(k)}/a_{kk}^{(k)}\)（\(i=k+1,\cdots,n\)），用 \(-m_{ik}\) 乘式（7.2.8）的第 \(k\) 个方程加上第 \(i\)（\(i=k+1,\cdots,n\)）个方程，消去第 \(k+1\) 个方程直到第 \(n\) 个方程的未知数 \(x_k\)，得到与式（7.2.1）等价的方程组 \(A^{(k+1)}x=b^{(k+1)}\)。

\(A^{(k+1)}\) 元素的计算公式为

\[
\begin{cases}
a_{ij}^{(k+1)}=a_{ij}^{(k)}-m_{ik}a_{kj}^{(k)}& (i,j=k+1,\cdots,n),\\
b_i^{(k+1)}=b_i^{(k)}-m_{ik}b_k^{(k)}& (i=k+1,\cdots,n).
\end{cases}
\tag{7.2.9}
\]

显然 \(A^{(k+1)}\) 的第 1 行直到第 \(k\) 行与 \(A^{(k)}\) 相同。

（3）继续这一过程，直到完成第 \(n-1\) 次消元。最后得到与原方程组等价的三角方程组

\[A^{(n)}x=b^{(n)}\tag{7.2.10}\]

\href{../../source-images/pdf-180.jpeg}{原页扫描}

或

\[
\begin{pmatrix}
a_{11}^{(1)}&a_{12}^{(1)}&\cdots&a_{1n}^{(1)}\\
&a_{22}^{(2)}&\cdots&a_{2n}^{(2)}\\
&&\ddots&\vdots\\
&&&a_{nn}^{(n)}
\end{pmatrix}
\begin{pmatrix}x_1\\x_2\\\vdots\\x_n\end{pmatrix}
=\begin{pmatrix}b_1^{(1)}\\b_2^{(2)}\\\vdots\\b_n^{(n)}\end{pmatrix}.
\]

由式（7.2.1）约化为式（7.2.10）的过程称为\textbf{消元过程}。

求解三角方程组（7.2.10），设 \(a_{ii}^{(i)}\ne0\)（\(i=1,2,\cdots,n-1\)），易得求解公式

\[
\begin{cases}
x_n=b_n^{(n)}/a_{nn}^{(n)},\\
x_k=\displaystyle\left(b_k^{(k)}-\sum_{j=k+1}^{n}a_{kj}^{(k)}x_j\right)/a_{kk}^{(k)}
\end{cases}
\quad(k=n-1,n-2,\cdots,2,1).
\tag{7.2.11}
\]

式（7.2.10）的求解过程称为\textbf{回代过程}。

如果 \(a_{11}^{(1)}=0\)，那么，由于 \(A\) 为非奇异矩阵，所以 \(A\) 的第 1 列一定有元素不等于零，例如 \(a_{i_1,1}\ne0\)，于是可交换两行元素（\(r_1\leftrightarrow r_{i_1}\)），将 \(a_{i_1,1}\) 调到第 1 行第 1 列的位置，然后进行消元计算，这时 \(A^{(2)}\) 右下角矩阵（\(n-1\) 阶）亦为非奇异矩阵。继续这一过程，Gauss 消去法照样可进行计算。

总结上述讨论即有如下定理。

\textbf{定理 7.1} 如果 \(A\) 为 \(n\) 阶非奇异矩阵，则可通过 Gauss 消去法（及交换两行的初等变换）将方程组（7.2.1）化为三角方程组（7.2.10）。

\(A\) 在什么条件下才能保证 \(a_{kk}^{(k)}\ne0\)（\(k=1,2,\cdots,n\)）？下面的引理给出了这个条件。

\textbf{引理} 约化的主元素 \(a_{ii}^{(i)}\ne0\)（\(i=1,2,\cdots,k\)）的充要条件是矩阵 \(A\) 的顺序主子式 \(D_i\ne0\)（\(i=1,2,\cdots,k\)），即

\[D_1=a_{11}\ne0,\]

\[
D_i=\begin{vmatrix}
a_{11}&\cdots&a_{1k}\\
\vdots&&\vdots\\
a_{i1}&\cdots&a_{ii}
\end{vmatrix}\ne0\quad(i=2,3,\cdots,k).
\tag{7.2.12}
\]

\textbf{证明} 利用归纳法证明引理的充分性。显然，当 \(k=1\) 时引理的充分性是成立的，现假设引理对 \(k-1\) 是成立的，求证引理对 \(k\) 亦成立。由归纳法，设 \(a_{ii}^{(i)}\ne0\)（\(i=1,2,\cdots,k-1\)），于是可用 Gauss 消去法将 \(A^{(1)}=A\) 约化到 \(A^{(k)}\) 中，即

\[
A^{(1)}\to A^{(k)}=\begin{pmatrix}
a_{11}^{(1)}&a_{12}^{(1)}&\cdots&\cdots&\cdots&a_{1n}^{(1)}\\
&a_{22}^{(2)}&\cdots&\cdots&\cdots&a_{2n}^{(2)}\\
&&\ddots&&&\vdots\\
&&&a_{kk}^{(k)}&\cdots&a_{kn}^{(k)}\\
&&&\vdots&&\vdots\\
&&&a_{nk}^{(k)}&\cdots&a_{nn}^{(k)}
\end{pmatrix},
\]

\begin{quote}
校注（agent 补充）：式（7.2.12）原书矩阵右上角清楚印作 \(a_{1k}\)，此处忠实保留。按同式左端 \(D_i\)、左下角 \(a_{i1}\) 和右下角 \(a_{ii}\)，第 \(i\) 阶顺序主子式的右上角应为 \(a_{1i}\)；这是疑似原书下标排印错误，不是 OCR 改写。参见\href{../assets/ch07a/pdf-180-principal-minor-detail.png}{原式放大裁图}。
\end{quote}

\href{../../source-images/pdf-181.jpeg}{原页扫描}

且有

\[
D_2=\begin{vmatrix}a_{11}^{(1)}&a_{12}^{(1)}\\0&a_{22}^{(2)}\end{vmatrix}
=a_{11}^{(1)}a_{22}^{(2)},\qquad
D_3=a_{11}^{(1)}a_{22}^{(2)}a_{33}^{(3)},
\]

\[
D_k=\begin{vmatrix}
a_{11}^{(1)}&a_{12}^{(1)}&\cdots&a_{1k}^{(1)}\\
&a_{22}^{(2)}&\cdots&a_{2k}^{(2)}\\
&&\ddots&\vdots\\
&&&a_{kk}^{(k)}
\end{vmatrix}
=a_{11}^{(1)}a_{22}^{(2)}\cdots a_{kk}^{(k)}.
\tag{7.2.13}
\]

由设 \(D_i\ne0\)（\(i=1,2,\cdots,k\)）及式（7.2.13），有 \(a_{kk}^{(k)}\ne0\)，即引理的充分性对 \(k\) 成立。

显然，由假设 \(a_{ii}^{(i)}\ne0\)（\(i=1,2,\cdots,k\)），利用式（7.2.13）亦可推出 \(D_i\ne0\)（\(i=1,2,\cdots,k\)）。

\textbf{推论} 如果 \(A\) 的顺序主子式 \(D_k\ne0\)（\(k=1,2,\cdots,n-1\)），则

\[
\begin{cases}
a_{11}^{(1)}=D_1,\\
a_{kk}^{(k)}=D_k/D_{k-1}\quad(k=2,3,\cdots,n).
\end{cases}
\]

\textbf{定理 7.2} 如果 \(n\) 阶矩阵 \(A\) 的所有顺序主子式均不为零，即 \(D_i\ne0\)（\(i=1,2,\cdots,n\)），则可通过 Gauss 消去法（不进行交换两行的初等变换），将方程组（7.2.1）约化为三角方程组（7.2.10）。

计算公式如下：

（1）消元计算（\(k=1,2,\cdots,n-1\)）。

\[m_{ik}=a_{ik}^{(k)}/a_{kk}^{(k)}\quad(i=k+1,\cdots,n),\]

\[a_{ij}^{(k+1)}=a_{ij}^{(k)}-m_{ik}a_{kj}^{(k)}\quad(i,j=k+1,\cdots,n),\]

\[b_i^{(k+1)}=b_i^{(k)}-m_{ik}b_k^{(k)}\quad(i=k+1,\cdots,n).\]

（2）回代计算。求解公式为式（7.2.11）。

\subsection{本节覆盖原页的校注记录}\label{ux672cux8282ux8986ux76d6ux539fux9875ux7684ux6821ux6ce8ux8bb0ux5f55}

下列记录按原页关联，含同页相邻小节的校注；计算时同时读取上文校注和完整原页。

\begin{itemize}
\tightlist
\item
  \texttt{ch07a-E001}：原书 PDF 页 180
\end{itemize}

\end{document}
